% ═══════════════════════════════════════════════════════════════
% MT-07b: Minitest Zahlen 21-100 (Números)
% Spanisch Klasse 7 - Sequenz 7b
% ═══════════════════════════════════════════════════════════════
% Dauer: 10-15 Minuten
% Punkte: 26 BE
% Inhalt: Zahlen 21-100, Preise, Einkaufsdialog
% ═══════════════════════════════════════════════════════════════

\documentclass[11pt, a4paper]{article}

\newif\ifswmodus
\swmodusfalse

\usepackage[utf8]{inputenc}
\usepackage[T1]{fontenc}
\usepackage[ngerman]{babel}
\usepackage{helvet}
\renewcommand{\familydefault}{\sfdefault}

\usepackage[margin=2cm]{geometry}
\usepackage{enumitem}
\usepackage{tabularx}
\usepackage{booktabs}

\usepackage{xcolor}
\usepackage{tcolorbox}
\tcbuselibrary{skins}

\usepackage{fancyhdr}
\usepackage{lastpage}
\usepackage{needspace}

% FARBEN
\definecolor{spanisch-color}{HTML}{c41e3a}
\definecolor{grau-color}{HTML}{888888}
\definecolor{hellgrau-color}{HTML}{f5f5f5}
\definecolor{orange-color}{HTML}{b35c00}

\definecolor{sw-dark}{HTML}{000000}
\definecolor{sw-gray}{HTML}{666666}
\definecolor{sw-white}{HTML}{ffffff}

\ifswmodus
    \colorlet{spanisch}{sw-dark}
    \colorlet{grau}{sw-gray}
    \colorlet{hellgrau}{sw-white}
    \colorlet{orange}{sw-dark}
\else
    \colorlet{spanisch}{spanisch-color}
    \colorlet{grau}{grau-color}
    \colorlet{hellgrau}{hellgrau-color}
    \colorlet{orange}{orange-color}
\fi

\newcommand{\fachfarbe}{spanisch}

\newcommand{\thema}{Minitest: Números 21-100}
\newcommand{\fach}{Spanisch}
\newcommand{\klasse}{7}
\newcommand{\maxpunkte}{26}
\newcommand{\zeit}{10-15 Minuten}
\newcommand{\datum}{\today}

\pagestyle{fancy}
\fancyhf{}
\renewcommand{\headrulewidth}{0.4pt}
\renewcommand{\footrulewidth}{0.4pt}

\lhead{\fach{} -- Klasse \klasse}
\chead{\textbf{MT-07b}}
\rhead{\datum}

\lfoot{\zeit}
\cfoot{}
\rfoot{Seite \thepage\ von \pageref{LastPage}}

\newcommand{\punktebox}[1]{\hfill\fbox{\small \_/#1}}
\newcommand{\luecke}[1]{\rule{#1}{0.4pt}}

\newtcolorbox{infobox}{
    colback=hellgrau,
    colframe=\fachfarbe,
    boxrule=1pt,
    arc=0pt,
    left=8pt, right=8pt, top=8pt, bottom=8pt
}

\newtcolorbox{hinweis}{
    colback=white,
    colframe=orange,
    boxrule=1pt,
    arc=0pt,
    left=5pt, right=5pt, top=5pt, bottom=5pt
}

\newenvironment{aufgabe}[1]{%
    \needspace{5\baselineskip}%
    \nopagebreak[4]%
    \vspace{10pt}
    \noindent\textbf{\textcolor{\fachfarbe}{Aufgabe #1}}
    \nopagebreak[4]%
    \vspace{5pt}
    \nopagebreak[4]%
    \begin{list}{}{\leftmargin=0pt}
    \item[]
}{%
    \end{list}
}

\newcommand{\es}[1]{\textcolor{\fachfarbe}{\textbf{#1}}}

\begin{document}

% ═══════════════════════════════════════════════════════════════
% KOPFBEREICH
% ═══════════════════════════════════════════════════════════════

\begin{center}
    {\Large\bfseries\textcolor{\fachfarbe}{\thema}}
\end{center}

\vspace{5pt}

\noindent
\begin{tabularx}{\textwidth}{|X|X|X|}
\hline
\textbf{Name:} \luecke{4cm} &
\textbf{Klasse:} \klasse &
\textbf{Datum:} \luecke{2cm} \\
\hline
\end{tabularx}

\vspace{10pt}

\begin{infobox}
\textbf{Zeit:} \zeit{} | \textbf{Punkte:} \maxpunkte{} BE | \textbf{Hilfsmittel:} keine
\end{infobox}

\vspace{15pt}

% ═══════════════════════════════════════════════════════════════
% AUFGABE 1: Zahlen als Wort schreiben (8 BE)
% ═══════════════════════════════════════════════════════════════

\begin{aufgabe}{1 -- Zahlen als Wort schreiben}
Schreibe die Zahl als spanisches Wort. \punktebox{8}
\end{aufgabe}

\noindent
a) 25 = \luecke{5cm}

\vspace{0.8em}

\noindent
b) 31 = \luecke{5cm}

\vspace{0.8em}

\noindent
c) 47 = \luecke{5cm}

\vspace{0.8em}

\noindent
d) 56 = \luecke{5cm}

\vspace{0.8em}

\noindent
e) 68 = \luecke{5cm}

\vspace{0.8em}

\noindent
f) 73 = \luecke{5cm}

\vspace{0.8em}

\noindent
g) 89 = \luecke{5cm}

\vspace{0.8em}

\noindent
h) 100 = \luecke{5cm}

\vspace{15pt}

% ═══════════════════════════════════════════════════════════════
% AUFGABE 2: Zahl zur Wort zuordnen (6 BE)
% ═══════════════════════════════════════════════════════════════

\begin{aufgabe}{2 -- Zahl zum Wort zuordnen}
Welche Zahl ist das? Schreibe die Zahl als Ziffer. \punktebox{6}
\end{aufgabe}

\noindent
a) \es{cuarenta y dos} = \luecke{2cm}

\vspace{0.8em}

\noindent
b) \es{setenta y cinco} = \luecke{2cm}

\vspace{0.8em}

\noindent
c) \es{noventa y nueve} = \luecke{2cm}

\vspace{0.8em}

\noindent
d) \es{veintiuno} = \luecke{2cm}

\vspace{0.8em}

\noindent
e) \es{sesenta y cuatro} = \luecke{2cm}

\vspace{0.8em}

\noindent
f) \es{ochenta y tres} = \luecke{2cm}

\vspace{15pt}

% ═══════════════════════════════════════════════════════════════
% AUFGABE 3: Preise verstehen (6 BE)
% ═══════════════════════════════════════════════════════════════

\begin{aufgabe}{3 -- Preise verstehen}
Schreibe den Preis als spanischen Satz. \punktebox{6}
\end{aufgabe}

\begin{hinweis}
\textbf{Beispiel:} 3,50 € = \es{Tres euros con cincuenta (céntimos).}
\end{hinweis}

\vspace{0.5em}

\noindent
a) 2,30 € = \luecke{8cm}

\vspace{1em}

\noindent
b) 5,75 € = \luecke{8cm}

\vspace{1em}

\noindent
c) 8,40 € = \luecke{8cm}

\vspace{15pt}

% ═══════════════════════════════════════════════════════════════
% AUFGABE 4: Einkaufsdialog (6 BE)
% ═══════════════════════════════════════════════════════════════

\begin{aufgabe}{4 -- Einkaufsdialog}
Ergänze den Dialog mit den passenden Fragen und Antworten. \punktebox{6}
\end{aufgabe}

\begin{hinweis}
\textbf{Hilfreiche Phrasen:} ¿Cuánto cuesta? | Cuesta ... euros. | ¿Cuánto es? | Son ... euros.
\end{hinweis}

\vspace{0.5em}

\noindent
\textbf{Kunde:} \luecke{6cm} \hfill (Wie viel kostet das?)

\vspace{0.8em}

\noindent
\textbf{Verkäufer:} \luecke{6cm} \hfill (Das kostet 4 Euro.)

\vspace{0.8em}

\noindent
\textbf{Kunde:} Quiero dos kilos. \luecke{4cm} \hfill (Wie viel macht das?)

\vspace{0.8em}

\noindent
\textbf{Verkäufer:} \luecke{6cm} \hfill (Das macht 8 Euro.)

% ═══════════════════════════════════════════════════════════════
% PUNKTEÜBERSICHT
% ═══════════════════════════════════════════════════════════════

\vspace{25pt}
\hrule
\vspace{10pt}

\noindent
\begin{tabularx}{\textwidth}{|X|c|c|c|c||c|}
\hline
\textbf{Aufgabe} & 1 & 2 & 3 & 4 & \textbf{Gesamt} \\
\hline
\textbf{max. BE} & 8 & 6 & 6 & 6 & \textbf{26} \\
\hline
\textbf{erreicht} & & & & & \\[0.8em]
\hline
\end{tabularx}

\vspace{10pt}

\begin{flushright}
\textbf{Note:} \luecke{2cm}
\end{flushright}

\end{document}
